\chapter{Product Management \& New-Product Development}\label{ch:product-management}

% Scope: product-market fit, discovery & prioritisation (RICE/Stage-Gate), roadmap and
%   build-vs-buy-vs-partner, new-product launch & adoption (Bass). NEW chapter (closes the PM/NPD gap).
% The running company IS a product; "what to build" is upstream of "what to charge" (\cref{ch:pricing}).
% Draft per house style: flowing prose + example/admonition set-offs; running-SaaS worked examples.

\section{Product-Market Fit}\label{sec:product-market-fit}
% complexity: 2

What does it mean for a product to fit its market? The question precedes every growth
conversation in this book. A company that scales acquisition before it has genuine
product-market fit is burning capital to fill a leaky bucket: customers arrive, they
do not stay, and the retention curve (\cref{eq:retention}) slopes sharply downward from
the first cohort. Product-market fit is the condition in which a product satisfies a
real, recurring need so completely that customers would be genuinely distressed to
lose access to it. Before that condition holds, increasing marketing spend accelerates
failure; after it holds, organic word-of-mouth starts contributing to growth, and the
economics of paid acquisition improve dramatically.

The canonical diagnostic was articulated by Sean Ellis after benchmarking nearly a
hundred early-stage technology companies. He asked a single, negatively framed survey
question: ``How would you feel if you could no longer use this product?'' The negative
framing strips away the politeness bias embedded in satisfaction scores. The threshold
is straightforward: if fewer than \qty{40}{\percent} of recently active users answer
``very disappointed,'' the product has not yet achieved fit with any segment it can
define. At \qty{40}{\percent} or above, empirical observation supports scaling.
\textcite{ellis2017hacking} documents the benchmarking work behind this threshold;
\textcite{cagan2018inspired} embeds it in a broader product-discovery framework.

The survey score is a leading indicator, not a definition. The true empirical proxy
for product-market fit is the shape of the cohort retention curve from \cref{eq:retention}.
A product with genuine fit displays an asymptotic survival curve: the curve flattens
toward a stable non-zero fraction rather than continuing to decline toward zero.
A curve that never flattens --- regardless of how many users report being very
disappointed --- means the product has a core-value problem that no amount of
growth investment will solve. The two diagnostics are complementary: the survey
catches problems early, while the retention curve confirms or refutes them with
behavioral evidence over time. \textcite{ries2011} articulates the build-measure-learn
logic that governs how quickly a team can iterate toward fit.

\begin{example}[PMF diagnostic for the running SaaS]
The running company's \qty{12}{\text{-month}} retention of \qty{92.5}{\percent} (\(c =
\qty{0.65}{\percent}/\text{mo}\)) is well above the empirical PMF bar: a \qty{92.5}{\percent}
annual survival rate produces an asymptotic curve, confirming that the product has genuine
fit with its current customer base. The concern is upstream: the trial-to-paid conversion
rate is only \qty{31}{\percent} (800 trials, 200 paid, per the funnel). This gap is not a
PMF failure --- activated, paying customers stay --- but an \emph{onboarding} failure.
Value is not surfacing fast enough during the trial. The diagnostic implication: instrument
the ``aha moment'' (in this product, running a first report), measure the time-to-first-report
for trial users, and test interventions that reduce it. High retention among paid users
confirms fit is real; low trial conversion confirms the path to fit is obscured.
\end{example}

\begin{admonition}{Warning}
Product-market fit is segment-specific, not product-wide. A \qty{40}{\percent} score
averaged across your entire user base can mask a \qty{65}{\percent} score for one segment
and a \qty{15}{\percent} score for another. Scaling into the wrong segment because the
blended score looked acceptable is how companies build large, high-churn customer bases
that look healthy until the first renewal cycle. Segment the survey by acquisition channel,
company size, and use-case before drawing conclusions. \textcite{torres2021continuous}
treats segment-specific discovery as the foundation of the continuous-discovery practice.
\end{admonition}

\section{Discovery \& Prioritisation}\label{sec:product-discovery}
% complexity: 4

Once a product has demonstrated fit with a defined segment, the question shifts from
\emph{does this product create value?} to \emph{which improvements create the most value,
fastest?} This is the prioritisation problem, and it is persistently distorted by
cognitive bias: loudest customer voice over median user need, engineer enthusiasm over
market impact, complexity of implementation over simplicity of benefit. The first
discipline of product management is replacing intuition about importance with a
structured estimate of expected value per unit of effort.

The RICE framework operationalises that estimate. Every proposed initiative is scored
on four dimensions, combined into a single comparable number:

\begin{equation}\label{eq:rice}
  \mathrm{RICE} \;=\; \frac{R \;\times\; I \;\times\; C}{E},
\end{equation}
where \(R\) is \emph{Reach} (the number of users or transactions affected per period),
\(I\) is \emph{Impact} (a tiered multiplier: 3 for massive, 2 for high, 1 for medium,
0.5 for low, 0.25 for minimal), \(C\) is \emph{Confidence} (probability-like certainty
in the Reach and Impact estimates: 1.0, 0.8, or 0.5), and \(E\) is \emph{Effort} in
person-months. The numerator rewards initiatives that reach many users with high
confidence of material impact; the denominator penalises large engineering commitments.
The formula's strength is forcing explicit estimates of all four factors, surfacing the
assumptions behind any gut-feel ranking.

The broader governance structure for more complex or hardware-intensive new-product
development is the Stage-Gate process (\cref{fig:stage-gate}). The framework segments
the development journey into discrete phases (Discovery, Scoping, Business Case,
Development, Testing, Launch), separated by decision gates where cross-functional
leadership makes an explicit Go, Kill, Hold, or Recycle call. Each gate prevents
investment from compounding on an initiative that has failed a necessary condition ---
market attractiveness at Gate 2, financial viability at Gate 3, technical feasibility
at Gate 4. \textcite{cooper2017winning} documents the empirical research behind the
process and its adaptations for software organisations; \textcite{cagan2018inspired}
describes how discovery and delivery teams run in parallel lanes rather than strict
sequence.

\begin{figure}[htbp]
  \centering
  \includegraphics[width=0.82\linewidth]{stage-gate}
  \caption{Stage-Gate process: five stages separated by Go/Kill/Hold decision gates.
    Diamonds represent gatekeeping decisions; rectangles represent stages of parallel
    cross-functional work. A ``Kill'' or ``Hold'' decision at any gate prevents further
    resource commitment until the blocking condition is resolved.}
  \label{fig:stage-gate}
\end{figure}

\begin{example}[RICE prioritisation: SSO vs.\ CSV export]
The product backlog contains two features. Feature A is single sign-on (SSO) integration,
primarily requested by enterprise prospects; Feature B is CSV data export, requested by
the operational users inside every existing account. Applying \cref{eq:rice} with
per-quarter estimates:

\medskip
\noindent Feature A (SSO): \(R = 200\) prospects, \(I = 3\) (high --- removes a procurement
blocker), \(C = 0.8\), \(E = 4\) person-months:
\[
  \mathrm{RICE}_A \;=\; \frac{200 \times 3 \times 0.8}{4} \;=\; 120.
\]
Feature B (CSV export): \(R = 500\) active users, \(I = 1\) (medium --- convenience
improvement), \(C = 0.9\), \(E = 1\) person-month:
\[
  \mathrm{RICE}_B \;=\; \frac{500 \times 1 \times 0.9}{1} \;=\; 450.
\]
B scores \num{3.75}\(\times\) higher. The implication: the bottleneck in the running
company's funnel is the \qty{31}{\percent} trial-to-paid conversion, not the enterprise
sales cycle (see \cref{sec:toc} for the constraint logic). Feature A would add friction-free
enterprise procurement --- valuable eventually --- but it does not address the constraint.
Feature B gives every current and prospective paid user a tangible workflow benefit, reaching
\num{500} users for a single person-month. Until the trial-to-paid rate is fixed, building
enterprise features is premature optimisation.
\end{example}

\begin{admonition}{Note}
RICE scores are not comparable across qualitatively different effort profiles.
A five-person-month architectural refactor and a one-person-month UI change sit in the
same denominator but carry incomparable risk, reversibility, and technical debt
implications. Use RICE to rank features within a roughly similar effort band; apply
Stage-Gate governance to initiatives large enough to require a formal business-case
gate. The framework suppresses bias; it does not replace judgment about what a
business needs. \textcite{torres2021continuous} cautions that frameworks applied
mechanically can themselves become a proxy for the thinking they were meant to enable.
\end{admonition}

\section{Roadmap \& Build-vs-Buy-vs-Partner}\label{sec:roadmap}
% complexity: 2

Even after RICE has ranked the backlog, a higher-order question arises for each major
initiative: should the capability be built in-house, bought (via licensing or acquisition),
or delivered through a partnership? This build-vs-buy-vs-partner decision shapes the
product roadmap more durably than any single feature choice, because it determines
whether the firm accumulates proprietary intellectual property or buys time at the cost
of strategic control. The financial criterion is \cref{eq:npv}: the option with the
higher net present value of expected margin contribution, properly charged for integration
and ongoing cost, is preferred. The strategic criterion is the VRIN test from
\cref{sec:rbv-vrin}: build only when the resulting capability is Valuable, Rare, Inimitable,
and Non-substitutable, because only then does the investment generate a durable moat.
Commodity functionality built in-house is opportunity-cost destruction.

The practical framework maps three routes against four dimensions: strategic control,
total cost of ownership, time to market, and vendor dependency. Building in-house
maximises control and IP ownership but carries the slowest delivery and highest execution
risk. Buying or licensing delivers the fastest time to market at the cost of vendor lock-in
and roadmap dependency. Partnering balances speed and capability with shared governance
and negotiated IP rights. The right choice is contingent on the capability's position in
the VRIN taxonomy, not on engineering preference.

\begin{example}[Integration marketplace: build vs.\ partner with Zapier]
The running company is considering how to serve the \qty{35}{\percent} of prospects who
ask for integrations with third-party tools (HR systems, CRMs, project trackers). Two
options: build a native integration marketplace (engineering cost \money{80000} upfront,
two engineers for four months, plus ongoing maintenance); or partner with an established
workflow automation platform at zero upfront cost in exchange for a \qty{30}{\percent}
revenue share on any plan upgrade attributed to integrations.

Applying \cref{eq:npv}: the native build requires \(I_0 = \money{80000}\) today. If
integrations reduce annual churn by \qty{1}{\percent} on the \money{4000000} ARR base,
the annual retention benefit is \money{40000}. Capitalised as a growing perpetuity at the
firm's \qty{11}{\percent} discount rate: \(\money{40000}/(0.11 - 0.03) = \money{500000}\).
NPV of building \(\approx \money{500000} - \money{80000} = \money{420000}\): clearly positive.

The partnership route saves \money{80000} upfront but pays a \qty{30}{\percent} rev-share
on attributed upgrades in perpetuity. More critically: the integration surface is what
locks customers into the product's data model and workflow. That is the definition of a
switching-cost moat (see \cref{sec:network-effects}). Outsourcing the integrations to a
third party means the stickiness accrues to the partner's platform, not the firm's.
VRIN logic confirms: a proprietary integration layer is Valuable (reduces churn),
Rare (few mid-market SaaS products build deeply), Inimitable (workflow knowledge is
embedded over time), and Non-substitutable. The build option wins --- not merely on NPV
but on strategic grounds.
\end{example}

\begin{admonition}{Warning}
Building commodity functionality --- features that any engineering team could replicate
in a weekend, that require no proprietary data, and that carry no network effect --- is
opportunity-cost destruction. Every person-month spent replicating what can be licensed
cheaply is a person-month not spent on the capability that constitutes the firm's
actual differentiation. The test is always VRIN, applied before the sprint plan is
written.
\end{admonition}

\section{New-Product Launch \& Adoption}\label{sec:adoption-curves}
% complexity: 4

Building the right product and pricing it correctly does not guarantee adoption. The
rate at which a market absorbs a new product is governed by social dynamics that
individual product decisions can influence but not override. The Bass Diffusion Model
is the foundational mathematical description of how those dynamics unfold. It
decomposes adoption into two behavioural forces: \emph{innovation} --- individuals
who adopt spontaneously, driven by external media and intrinsic affinity for new
technology --- and \emph{imitation} --- individuals who adopt because they observe
peers who already have. The ratio of these two forces determines not only the eventual
adoption level but the timing and height of the peak.

In the discrete-time form used for practical forecasting, new adopters at period \(t\)
are:
\begin{equation}\label{eq:bass}
  N(t) \;=\; \bigl[p + q\,\tfrac{N(t-1)}{M}\bigr]\,\bigl[M - N(t-1)\bigr],
\end{equation}
where \(M\) is the total addressable market (ultimate adopter ceiling), \(p\) is the
coefficient of innovation, \(q\) is the coefficient of imitation, and \(N(t-1)\) is
the cumulative installed base at the start of period \(t\). The first bracketed term
is the instantaneous adoption hazard: the probability that a non-adopter converts,
combining a base rate (\(p\)) with a social-pressure term that rises as the installed
base grows (\(q \cdot N(t-1)/M\)). The second bracketed term is the pool of remaining
non-adopters. When the installed base is small, the social term is negligible and
growth is slow; as the base builds, the social term accelerates adoption until market
saturation compresses the remaining pool and growth decelerates. The result is the
characteristic S-shaped cumulative adoption curve and bell-shaped new-adopter curve
shown in \cref{fig:bass-diffusion}. Meta-analytic estimates across diverse product
categories set \(\bar{p} \approx 0.03\) and \(\bar{q} \approx 0.38\)
\parencite{bass1969newproduct}.

\begin{figure}[htbp]
  \centering
  \includegraphics[width=0.86\linewidth]{bass-diffusion}
  \caption{Bass diffusion model: cumulative adoption S-curves for the enterprise-tier
    launch (\(M = 500\), \(p = 0.02\), \(q = 0.3\)). The \emph{innovator-only} curve
    (coefficient \(p\) acting alone, no imitation) tracks purely linear diffusion; the
    full Bass curve bends into a sigmoid once the installed base is large enough for
    word-of-mouth to dominate. The dashed vertical line marks the new-adopter peak
    (month 10, \(\approx\num{43}\) new accounts).}
  \label{fig:bass-diffusion}
\end{figure}

\begin{example}[Enterprise-tier launch: Bass calibration]
The running company launches an enterprise pricing tier targeting \(M = 500\) mid-market
accounts in its addressable segment. Historical outbound prospecting suggests a base
conversion rate of \(p = 0.02\) (innovators --- accounts that evaluate and buy
independently). Customer success interviews suggest word-of-mouth referrals within
industry peer networks are active, estimated at \(q = 0.3\).

Applying \cref{eq:bass} iteratively:

\medskip
\textbf{Month 1.}\ No installed base yet. New adopters \(= (0.02 + 0.3 \times 0/500)
\times (500 - 0) = 0.02 \times 500 = 10\). Cumulative: \num{10}.

\textbf{Month 2.}\ Social term now contributes. New adopters \(\approx (0.02 + 0.3
\times 10/500) \times 490 = 0.026 \times 490 \approx 13\). Cumulative: \num{23}.

The acceleration continues through month 10, where the new-adopter count peaks at
approximately \num{43} per month (cumulative \(\approx\num{268}\) of \num{500}).
After the peak the remaining addressable pool shrinks faster than the social term
grows, and new-adopter counts decline back toward the innovator-only floor.

The managerial implication: the launch plan should front-load outreach in months 1--3,
securing the early innovators whose presence seeds the social term. Each early customer
reference --- case study, peer referral, industry conference testimony --- raises the
effective \(q\) and advances the peak. This is precisely the role of the minimum
viable product in the lean-startup cycle (\cref{sec:lean-startup}): an MVP placed
with five to ten early adopters calibrates both \(p\) and \(q\) empirically before
full investment in the launch funnel.
\end{example}

\begin{admonition}{Warning}
The Bass model describes macro-level diffusion across a population of independent
potential adopters --- it is not a cohort CLV model. \Cref{eq:bass} says nothing about
what each individual account is worth after adoption, or about the \emph{retention}
dynamics within the installed base once acquired. The two models are complementary:
Bass forecasts the supply of new accounts over time; \cref{eq:retention} and \cref{eq:clv}
in \cref{ch:customer-economics} determine the value extracted from each account after
it joins. Confusing the two --- using Bass adoption curves to project revenue without
a separate cohort-retention analysis --- is the most common modelling error in
launch-stage financial planning. \textcite{bass1969newproduct} is explicit on scope:
the model governs first-purchase diffusion, not repeat purchase or retention.
\end{admonition}
